% Generado por sistema/backend ReporteTex.java -- no editar a mano.
% Datos: Base de datos: 148 importado (titulados) / Base de datos: 53 importado (empleadores)
\subsubsection*{C3. Formación complementaria $\times$ remuneración}

\textit{Propósito.} Determinar si la formación complementaria se asocia con el nivel de remuneración, como sustento para la oferta de actualización profesional y posgrado. Respalda el criterio ARCU-SUR 2.4.1 / 3.2.3. Se aplica sobre: trabajan en una organización, con nivel de significancia $\alpha = 0.05$.

\textit{Paso 1. Contar con la muestra: tabla de contingencia} (frecuencia observada $FO_{ij}$, $N = 108$).

\begin{mitabla}{C3: frecuencia observada $FO_{ij}$}
\begin{tabular}{lcc|c}
\toprule
Formación complementaria & Hasta Bs 5.000 & Más de Bs 5.000 & Total fila \\
\midrule
Sí & 21 & 52 & 73 \\
No & 13 & 22 & 35 \\
\midrule
Total columna & 34 & 74 & $N=108$ \\
\bottomrule
\end{tabular}
\end{mitabla}

\textit{Paso 2. Calcular la frecuencia esperada} de cada celda, $FE_{ij} = (\text{total fila } i)(\text{total columna } j)/N$. Por ejemplo, para la celda ``Sí -- Hasta Bs 5.000'': $FE_{11} = 73 \times 34 / 108 = 22.9815$.

\begin{mitabla}{C3: frecuencia esperada $FE_{ij}$}
\begin{tabular}{lcc}
\toprule
Formación complementaria & Hasta Bs 5.000 & Más de Bs 5.000 \\
\midrule
Sí & $73\cdot34/108 = 22.98$ & $73\cdot74/108 = 50.02$ \\
No & $35\cdot34/108 = 11.02$ & $35\cdot74/108 = 23.98$ \\
\bottomrule
\end{tabular}
\end{mitabla}

\textit{Paso 3. Calcular el estadístico chi-cuadrado,} sumando $(FO_{ij}-FE_{ij})^2/FE_{ij}$ de cada celda.

\begin{mitabla}{C3: cálculo de $X_0^2$ celda por celda}
\begin{tabular}{L{5.4cm}cccc}
\toprule
Celda & $FO_{ij}$ & $FE_{ij}$ & $(FO_{ij}-FE_{ij})^2$ & $(FO_{ij}-FE_{ij})^2/FE_{ij}$ \\
\midrule
Sí -- Hasta Bs 5.000 & 21 & 22.9815 & 3.9263 & 0.1708 \\
Sí -- Más de Bs 5.000 & 52 & 50.0185 & 3.9263 & 0.0785 \\
No -- Hasta Bs 5.000 & 13 & 11.0185 & 3.9263 & 0.3563 \\
No -- Más de Bs 5.000 & 22 & 23.9815 & 3.9263 & 0.1637 \\
\midrule
\multicolumn{4}{r}{$X_0^2 = \sum$} & 0.7694 \\
\bottomrule
\end{tabular}
\end{mitabla}

\textit{Paso 4. Determinar el estadístico teórico.} Con $k = 2$ filas y $m = 2$ columnas, los grados de libertad son $(k-1)(m-1) = (2-1)(2-1) = 1$. De la tabla de chi-cuadrada: $X_{0.05,\,1}^2 = 3.841$.

\textit{Paso 5. Establecer la relación entre estadísticos:} $X_0^2 < X_{0.05,\,1}^2 \;\Rightarrow\; 0.7694 < 3.841$.

\textit{Paso 6. Conclusión.}

\begin{center}
\begin{tabular}{l}
Si ($X_0^2 < X_{0.05,\,1}^2$)\\
\quad \fbox{Independencia}\\
Sino\\
\quad Relación entre las variables\\
FinSi
\end{tabular}
\end{center}

No se puede rechazar que formación complementaria y remuneración sean independientes.

